\section{Prompts and Templates}
\label{sec:app_prompts}
% =============================================================================

Three prompt regimes govern the pipeline: Stage~1 projector alignment, Stage~2 DAPO
rollouts, and the shared \texttt{<invoke>} tool-call format.

%------------------------------------------------------------------------------
\subsection{Stage~1: Projector Alignment}
\label{sec:app_prompts_stage1}

Stage~1 trains the LatentBridge projector ($H_\varphi$, the linear adapter mapping BT4
residual activations into the LLM token space) via supervised learning on chess
instruction data.
Each episode presents several independent questions about the same board position; the
\texttt{<state>} placeholder marks the $k{=}32$ state tokens projected from the BT4
residual stream and inserted into the LLM's context at that point.
Four question types are sampled per FEN: position evaluation, principal variation, legal
moves, and brief verbal description.
Move notation alternates UCI and SAN with probability 70\,/\,30\,\%; evaluations are in
pawn units.

\begin{tcolorbox}[promptbox,
  title={$\triangleright$\; Stage~1 System Prompt}]
\begin{quote}
You are LLAMIA, a chess reasoning assistant.
The board state is encoded in the \texttt{<state>} token sequence embedded in the
user's message.
Answer the question about this position precisely and concisely.
Do not narrate your reasoning; give the answer directly.
\end{quote}
\end{tcolorbox}

\begin{tcolorbox}[notebox, title={Stage~1 Training Format}]
\small
\textbf{User:}~What is the evaluation of the following position?
Describe the evaluation followed by the continuation in UCI format, e.g.\
\emph{``Evaluation is +1.5, Advantage for White, e2e4 e7e5''}. \texttt{<state>}\\[3pt]
\textbf{LLAMIA:}~Evaluation is $+0.5$, white has a slight advantage,
\texttt{d5c6 b4c4 e4d6 c4d3 \ldots}

\tcbline
\textbf{User:}~Give me the principal variation(s); just the evaluation(s) and the
line(s). \texttt{<state>}\\[3pt]
\textbf{LLAMIA:}~$+0.4$, \texttt{e4d6 a3a4 d5c6 b4c3 c6d5 c3d3}

\tcbline
\textbf{User:}~What are the legal moves in this position?
Give me up to 10 comma-separated UCI moves and nothing else. \texttt{<state>}\\[3pt]
\textbf{LLAMIA:}~\texttt{e4f2, d5e6, e4c3, e4g3, e4d6, e4g5, e4d2, e4f6, d5d6, e4c5}

\tcbline
\textbf{User:}~Describe the current position briefly. \texttt{<state>}\\[3pt]
\textbf{LLAMIA:}~It is black's turn. White's King is somewhat exposed, with threats from
Black's Knight and Pawn. Black has a centralised active Knight and Pawn. This position
appears as a complex endgame with White holding a slight advantage due to the extra pawn.
\end{tcolorbox}

%------------------------------------------------------------------------------
\subsection{Stage~2: DAPO System Prompt}
\label{sec:app_prompts_stage2}

\paragraph{Toy task (puzzle popularity / Elo).}
The system prompt below is used verbatim during DAPO rollouts for the toy task
(\S\ref{sec:app_stage2}).
The explicit refusal-suppression clause is required: without it, Qwen3-4B defaults to
``I cannot directly determine\ldots'' and never emits a tool call, collapsing the
format-pass rate to 0\,\% (verified on 20 held-out puzzles at greedy decoding).

\begin{tcolorbox}[promptbox,
  title={$\triangleright$\; Stage~2 System Prompt --- Puzzle Understanding}]
\begin{quote}
You are a chess expert with access to lc0, a top neural network engine.
A puzzle position has been loaded from the FEN in the user's message.
Use the tools below to analyse the position, then estimate the puzzle's popularity and
difficulty rating.

\smallskip
\textbf{Available tools:}
\begin{itemize}[noitemsep,topsep=1pt,leftmargin=14pt]
  \item \texttt{get\_position} --- FEN, ASCII board, legal moves.
  \item \texttt{analyze(nodes, multipv)} --- lc0 MCTS search.
  \item \texttt{get\_policy(nodes)} --- raw NN priors and per-move values.
\end{itemize}

\smallskip
\textbf{You must always provide a numeric answer.}
Refusing or writing ``I cannot determine'' is not permitted ---
give your best guess even if uncertain.

\smallskip
End your reply with exactly one line in the form:\\[2pt]
\texttt{The popularity is \textit{<int>} and the ELO is \textit{<int>}}
\end{quote}
\end{tcolorbox}

\paragraph{Full LLAMIA-Bench.}
All four task families (behaviour cloning, puzzle understanding, move
annotation, game commentary) share the same tool catalogue as the toy task.
Task-specific instructions replace the popularity/Elo mandate; the
refusal-suppression clause is retained across all variants.

%------------------------------------------------------------------------------
\subsection{\texttt{<invoke>} Trigger Format}
\label{sec:app_prompts_invoke}

Latent invocation (\texttt{<invoke>}) fires after every board-mutating tool call
(\texttt{make\_move}, \texttt{undo\_move}, \texttt{reset\_position}).
The harness re-runs the BT4 forward pass on the updated FEN, encodes a fresh set of
$k{=}32$ state tokens, and prepends them as a \texttt{<state>} prefix to the next user
turn.
Because mutating calls are dispatched sequentially (\S\ref{sec:app_toolcall}), the
re-encoding always sees a fully settled board state.
Read-only calls (\texttt{analyze}, \texttt{get\_policy}, \texttt{get\_position}) do not
trigger re-encoding; the token budget is therefore capped at $k$ additional tokens per
state transition, regardless of analysis depth.

\begin{tcolorbox}[notebox, title={\texttt{<invoke>} Turn Structure: Blunder Analysis}]
\small
\textbf{Round~1 --- User:} \texttt{<state$_0$>} Why is \texttt{Nxd4} a blunder here?
[FEN $s_0$]\\[3pt]
\textbf{Round~1 --- LLAMIA:} \textit{(emits two parallel \texttt{analyze} calls ---
see \S\ref{sec:app_toolcall_example})}

\tcbline
\textbf{Round~2 --- Tool:} \textit{(harness returns engine evaluations; state unchanged)}\\[3pt]
\textbf{Round~2 --- LLAMIA:} Final answer citing \texttt{Qa5+} / \texttt{Qxg5}.
\end{tcolorbox}

Since no board mutation occurs in blunder analysis, \texttt{<invoke>} does not fire.
The re-encoding path is active in multi-step planning episodes, where the agent sequences
\texttt{make\_move} calls to explore a variation before deciding on a recommendation.

\begin{tcolorbox}[notebox, title={\texttt{<invoke>} Turn Structure: Multi-step Planning}]
\small
\textbf{Round~1 --- User:} \texttt{<state$_0$>} Find the best three-move combination.
[FEN $s_0$]\\[3pt]
\textbf{Round~1 --- LLAMIA:} \texttt{make\_move(e2e4)}

\tcbline
\textbf{Round~2 --- Tool:} \{move\_played: \texttt{e4}, fen: $s_1$, \ldots\}\\
\textit{(harness re-encodes $s_1$ $\!\to\!$ fresh \texttt{<state$_1$>})}\\[3pt]
\textbf{Round~2 --- User (injected):} \texttt{<state$_1$>}\\[3pt]
\textbf{Round~2 --- LLAMIA:} \texttt{analyze(multipv=3)}

\tcbline
\textbf{Round~3 --- Tool:} \textit{(engine lines from $s_1$)}\\[3pt]
\textbf{Round~3 --- LLAMIA:} \texttt{undo\_move()}

\tcbline
\textbf{Round~4 --- Tool:} \{fen: $s_0$, \ldots\}\\
\textit{(harness re-encodes $s_0$ $\!\to\!$ fresh \texttt{<state$_0'$>})}\\[3pt]
\textbf{Round~4 --- User (injected):} \texttt{<state$_0'$>}\\[3pt]
\textbf{Round~4 --- LLAMIA:} Final recommendation.
\end{tcolorbox}

\noindent The injected \texttt{<state>} prefix in Rounds~2 and~4 is invisible to the
human user; the harness inserts it programmatically before forwarding the tool result to
the next LLM call, keeping the re-encoding fully transparent to the model's reasoning loop.

%------------------------------------------------------------------------------
\subsection{Inference Harness and Tool-Call Protocol}
\label{sec:app_toolcall}

The inference harness connects the LLM to lc0 via a six-tool stateful API.
Read-only calls (\texttt{analyze}, \texttt{get\_policy}, \texttt{get\_position}) are
dispatched in parallel; mutating calls (\texttt{make\_move}, \texttt{undo\_move},
\texttt{reset\_position}) are dispatched sequentially to preserve board consistency.

\subsubsection{Agent State}
\label{sec:app_toolcall_state}
\begin{tcolorbox}[statebox, title={Agent State}]
The agent maintains two fields across tool calls: the current \textbf{board position}
(FEN, castling rights, en-passant square, 50-move clock) and a \textbf{move history}
(SAN list from episode start).
Derived on demand: ASCII board, turn, move number, in-check flag, and legal move list
(truncated to 24 entries; full count reported via \texttt{total\_legal}).
Move notation is accepted as UCI (\texttt{e2e4}) or SAN (\texttt{Nf3}, \texttt{O-O});
illegal moves return \texttt{\{"error": \ldots\}} without raising, allowing the LLM to
retry with a corrected move.
State resets at episode start or on an explicit \texttt{reset\_position} call.
\end{tcolorbox}

\subsubsection{System Prompt}
\label{sec:app_toolcall_prompt}

\begin{tcolorbox}[promptbox, title={$\triangleright$\; System Prompt --- Chess Analyst}]
\begin{quote}
You are a chess expert with access to lc0, a top neural network engine.
A board is already loaded---any FEN in the user's query has been applied for you.
Do not call \texttt{reset\_position} unless you need a different position.

\smallskip
Tools (the board is stateful across calls):

\{\{ \$TOOL DEFINITIONS\}\}
\end{quote}
\end{tcolorbox}

\subsubsection{Tool Definitions}
\label{sec:app_toolcall_tools}
\label{sec:app_toolcall_example}

\begin{tcolorbox}[toolbox, title={Tool API Summary}]
\renewcommand{\arraystretch}{1.2}
\begin{adjustbox}{max width=\linewidth}
\begin{tabular}{@{}p{4.5cm}p{4.2cm}p{6cm}@{}}
\toprule
\textbf{Tool} & \textbf{Parameters} & \textbf{Returns} \\
\midrule
\texttt{get\_position()} & --- & FEN, ASCII board, turn, in-check, legal moves ($\leq$24), history \\
\texttt{make\_move(move)} & \texttt{move}: UCI or SAN & move played, FEN, turn, in-check, checkmate \\
\texttt{undo\_move()} & --- & FEN, turn \\
\texttt{reset\_position(fen)} & \texttt{fen}: FEN string & FEN, status; locked if FEN was auto-loaded \\
\texttt{analyze(nodes, multipv, moves)} & defaults: 800, 3, \texttt{[]} & best move, PV lines in SAN with $100\!\cdot\!Q$ scores \\
\texttt{get\_policy(nodes)} & default: auto ($\#$legal$+2$) & per-move $P$ (prior), $V$ (value), $Q$ (action-value), $N$ (visits) \\
\bottomrule
\end{tabular}
\end{adjustbox}
\smallskip

\noindent PVs from \texttt{analyze} are translated from UCI to SAN by the harness.
\texttt{get\_policy} runs the minimum search for a single visit per root child and returns
raw NN beliefs before MCTS modifies them.

\smallskip
\noindent\textbf{Note (LLAMIA only):} Calling \texttt{get\_policy} causes the harness to run a fresh BT4 forward pass on the current position
and inject $k{=}32$ latent state tokens into the next LLM turn, in addition to the text return.
In LLAMIA-Verb, only the text is returned.
\end{tcolorbox}

% =============================================================================
